\documentclass[final,t]{beamer}
\mode<presentation>
{
  \usetheme{NeSI}
%  \usetheme{Aachen}
%  \usetheme{Oldi6}
%  \usetheme{I6td}
%  \usetheme{I6dv}
%  \usetheme{I6pd}
%  \usetheme{I6pd2}
}
% additional settings
\setbeamerfont{itemize}{size=\normalsize}
\setbeamerfont{itemize/enumerate body}{size=\normalsize}
\setbeamerfont{itemize/enumerate subbody}{size=\normalsize}

% additional packages
\usepackage{times}
\usepackage{amsmath, amsthm, amssymb, soul, color, multicol, type1cm, verbatim, latexsym, float,multicol}
\usepackage{exscale}
%\boldmath
\usepackage{booktabs, array}
%\usepackage{rotating} %sideways environment
\usepackage[english]{babel}
\usepackage[latin1]{inputenc}
%\usepackage[orientation=landscape,size=custom,width=200,height=120,scale=1.9]{beamerposter}
\usepackage[orientation=landscape,size=a4]{beamerposter}
\listfiles
\graphicspath{{figures/}}
% Display a grid to help align images
%\beamertemplategridbackground[1cm]

\usepackage{listings}
\usepackage{xspace}
\usepackage{fp}
\usepackage{ifthen}

\title{\huge High Performance Computing Quick Reference Guide to the NeSI Pan Cluster}
\author{New Zealand eScience Infrastructure (\href{https://www.nesi.org.nz/}{NeSI}) @ \href{http://www.eresearch.auckland.ac.nz/}{The University of Auckland}}
\institute[]{
NeSI is a partnership for all New Zealand researchers delivered by 
  \href{http://www.eresearch.auckland.ac.nz/}{The University of Auckland}, \href{http://www.niwa.co.nz/our-services/hpcf}{NIWA}, \href{http://www.bluefern.canterbury.ac.nz/}{University of Canterbury}, \\
  \href{http://www.landcareresearch.co.nz/}{Landcare Research}, \href{http://www.otago.ac.nz/}{Otago University} and \href{http://msi.govt.nz/}{MBIE}
}

% abbreviations
\makeatletter
\makeatother

%%%%%%%%%%%%%%%%%%%%%%%%%%%%%%%%%%%%%%%%%%%%%%%%%%%%%%%%%%%%%%%%%%%%%%%%%%%%%%%%%%%%%%%%%%%%%%%%%%%%%%%%%%%%
%%%%%%%%%%%%%%%%%%%%%%%%%%%%%%%%%%%%%%%%%%%%%%%%%%%%%%%%%%%%%%%%%%%%%%%%%%%%%%%%%%%%%%%%%%%%%%%%%%%%%%%%%%%%
\begin{document}
\begin{frame}[fragile]{} 
  \begin{columns}[t]
    \begin{column}{.32\linewidth}

      %%%%%%%%%%%%%%%%%%%%%%%%%%%%%%%%%%%%%%%%%%%%%%%%%%%%%%%%%%%%%%%%%%%%%%%%%%%%%%%%%%%%%%%%%%%%%%%%%%%%%%%%%%%%
      \begin{block}{Contact Information}
        \begin{itemize}
        \item \textbf{web}: \url{http://www.nesi.org.nz}
        \item \textbf{wiki}: \url{https://wiki.auckland.ac.nz/display/CER/}
        \item \textbf{ganglia}: \url{http://ganglia.uoa.nesi.org.nz}
        \item \textbf{support} : \href{mailto:support@nesi.org.nz}{support@nesi.org.nz}
        \end{itemize}
        \vspace*{-1.5cm}
        \hspace*{7cm}
        \includegraphics[width=1.4cm]{img/qrcode.png} 
      \end{block}
      %%%%%%%%%%%%%%%%%%%%%%%%%%%%%%%%%%%%%%%%%%%%%%%%%%%%%%%%%%%%%%%%%%%%%%%%%%%%%%%%%%%%%%%%%%%%%%%%%%%%%%%%%%%%

      %%%%%%%%%%%%%%%%%%%%%%%%%%%%%%%%%%%%%%%%%%%%%%%%%%%%%%%%%%%%%%%%%%%%%%%%%%%%%%%%%%%%%%%%%%%%%%%%%%%%%%%%%%%%
      \begin{block}{Recommended Best Practices}
      \begin{itemize}
         \item Use the \textbf{SubmitScripts Templates} that are in /share/SubmitScripts.
         \item In order to improve the scalability of your MPI code, we suggest to minimise the number of nodes to be used. The \verb|--nodes=N-M| Slurm option, will define the minimum (N) and the maximum (M) number of nodes.
        \end{itemize}
      \end{block}
      %%%%%%%%%%%%%%%%%%%%%%%%%%%%%%%%%%%%%%%%%%%%%%%%%%%%%%%%%%%%%%%%%%%%%%%%%%%%%%%%%%%%%%%%%%%%%%%%%%%%%%%%%%%%

      %%%%%%%%%%%%%%%%%%%%%%%%%%%%%%%%%%%%%%%%%%%%%%%%%%%%%%%%%%%%%%%%%%%%%%%%%%%%%%%%%%%%%%%%%%%%%%%%%%%%%%%%%%%%
      \begin{block}{Useful Quick References}
        \vspace*{-3ex}
        \begin{multicols}{2}
        \begin{itemize}
        \item \href{https://wiki.auckland.ac.nz/download/attachments/37950883/vi+Quick+Reference.pdf}{VI Quick Reference}
        \item \href{https://wiki.auckland.ac.nz/download/attachments/37950883/Bash+Quick+Reference.pdf}{BASH Quick Reference} 
        \item \href{https://wiki.auckland.ac.nz/download/attachments/65014423/fwunixref.pdf?version=1&modificationDate=1399584680387&api=v2}{Linux Quick Reference}
        \item \href{https://wiki.auckland.ac.nz/download/attachments/57639811/OpenMP3.1-FortranCard.pdf?version=1&modificationDate=1399506093589&api=v2}{OpenMP Fortran Syntax}
        \item \href{https://wiki.auckland.ac.nz/download/attachments/57639811/OpenMP3.1-CCard.pdf?version=1&modificationDate=1399506057930&api=v2}{OpenMP 3.1 API C/C++ Syntax}
        \item \href{https://wiki.auckland.ac.nz/download/attachments/57639811/mpi-quick-ref.pdf?version=1&modificationDate=1399506030365&api=v2}{MPI Quick Reference}
        \end{itemize}
        \end{multicols}
        \vspace*{-3ex}
      \end{block}
      %%%%%%%%%%%%%%%%%%%%%%%%%%%%%%%%%%%%%%%%%%%%%%%%%%%%%%%%%%%%%%%%%%%%%%%%%%%%%%%%%%%%%%%%%%%%%%%%%%%%%%%%%%%%

      %%%%%%%%%%%%%%%%%%%%%%%%%%%%%%%%%%%%%%%%%%%%%%%%%%%%%%%%%%%%%%%%%%%%%%%%%%%%%%%%%%%%%%%%%%%%%%%%%%%%%%%%%%%%
      \begin{block}{NeSI Pan Cluster}
        \vspace*{-2ex}
      \begin{center}
      \begin{tabular}{|c|c|c|c|}
      \hline 
      \textbf{Architecture} & \textbf{Westmere} & \textbf{SandyBridge} & \textbf{LargeMem} \\ 
      \hline 
      Model & X5660 & E5-2680 & E7-4870 \\ 
      \hline 
      Clock Speed & 2.8 GHz & 2.7 GHz & 2.4GHz \\ 
      \hline 
      Cache & 12MB & 20MB & 30MB \\ 
      \hline 
      Intel QPI speed & 6.4GT/s & 8 GT/s & 6.4GT/ \\ 
      \hline 
      Cores/socket & 6 & 8 & 10 \\ 
      \hline       
      Cores/node & 12 & 16 & 40 \\ 
      \hline 
      Mem/node & 96GB & 128GB & 512GB \\ 
      \hline 
      GFLOPS/node & 134.4 & 345.6 & 384.0 \\ 
      \hline 
      \# nodes & 76 & 194 & 4 \\ 
      \hline 
      \end{tabular} 
      \end{center}
        \vspace*{-2ex}
      \end{block}
      %%%%%%%%%%%%%%%%%%%%%%%%%%%%%%%%%%%%%%%%%%%%%%%%%%%%%%%%%%%%%%%%%%%%%%%%%%%%%%%%%%%%%%%%%%%%%%%%%%%%%%%%%%%%


      %%%%%%%%%%%%%%%%%%%%%%%%%%%%%%%%%%%%%%%%%%%%%%%%%%%%%%%%%%%%%%%%%%%%%%%%%%%%%%%%%%%%%%%%%%%%%%%%%%%%%%%%%%%%
      \begin{block}{NeSI Pan Cluster - Co-Processors}
        \vspace*{-2ex}
      \begin{center}
      \begin{tabular}{|c|c|c|c|}
      \hline 
      \textbf{Architecture} & \textbf{Nvidia Fermi} & \textbf{Nvidia Kepler} & \textbf{Intel Phi} \\ 
      \hline 
      Main CPU                  &  X5660/E5-2680       &  E5-2680                   &  E5-2680 \\
      \hline 
      Model                       &  M2090                     &   K20X                       &  5110P \\ 
      \hline 
      Clock Speed             &  1.3GHz                    &  0.732GHz                      &  1.053GHz\\ 
      \hline 
      Cores/Dev.             & 512                            & 2688                           & 60 (240) \\ 
      \hline       
      Dev./node              &  2                             &   2                                &  2  \\ 
      \hline 
      Mem/Dev.             & 6GB                          & 6GB                               & 8GB  \\ 
      \hline 
      TFLOPS/Dev         & 1.33                          & 1.17                              & 1.01 \\ 
      \hline 
      \# nodes               & 16                             & 5                                 & 2 \\ 
      \hline 
      \end{tabular} 
      \end{center}
        \vspace*{-2ex}
      \end{block}
      %%%%%%%%%%%%%%%%%%%%%%%%%%%%%%%%%%%%%%%%%%%%%%%%%%%%%%%%%%%%%%%%%%%%%%%%%%%%%%%%%%%%%%%%%%%%%%%%%%%%%%%%%%%%


      %%%%%%%%%%%%%%%%%%%%%%%%%%%%%%%%%%%%%%%%%%%%%%%%%%%%%%%%%%%%%%%%%%%%%%%%%%%%%%%%%%%%%%%%%%%%%%%%%%%%%%%%%%%%
      \begin{block}{Disk Spaces + Default Quota}
        \vspace*{-2ex}
      \begin{center}
      \begin{tabular}{|l|c|c|c|c|c|c|}
      \hline 
      \textbf{FileSystem} & \textbf{Space} & \textbf{Quota} & \textbf{ACL} & \textbf{Backup} & \textbf{Type} & \textbf{Usage}  \\ 
      \hline 
      \$HOME & 120TB & 2GB & rw & yes & GPFS & Archive \\ 
      \hline 
      /project & 120TB & 30GB & rw & yes & GPFS & Archive \\ 
      \hline 
      /share & 120TB & - & ro & yes & GPFS & Archive \\ 
      \hline 
      \$TMP\_DIR & 240GB & - & rw & NO & EXT4 & io \\ 
      \hline 
      \$SCRATCH\_DIR & 13TB & - & rw & NO & GPFS & io \\ 
      \hline 
      \$SHM\_DIR & Mem & - & rw & NO & RamFS & io \\ 
      \hline 
      \end{tabular} 
      \end{center}
        \vspace*{-2ex}
      \end{block}
      %%%%%%%%%%%%%%%%%%%%%%%%%%%%%%%%%%%%%%%%%%%%%%%%%%%%%%%%%%%%%%%%%%%%%%%%%%%%%%%%%%%%%%%%%%%%%%%%%%%%%%%%%%%%
 
 


 

     \end{column}

%%%%%%%%%%%%%%%%%%%%%%%%%%%%%%%%%%%%%%%%%%%%%%%%%%%%%%%%%%%%%%
%%%%%%%%%%%%%%%%%%%%%%%%%%%%%%%%%%%%%%%%%%%%%%%%%%%%%%%%%%%%%%
%%%%%%%%%%%%%%%%%%%%%%%%%%%%%%%%%%%%%%%%%%%%%%%%%%%%%%%%%%%%%%
    
    \begin{column}{.32\linewidth}
    \vskip-2ex
    \begin{columns}[t]
    \begin{column}{.54\linewidth}
      %%%%%%%%%%%%%%%%%%%%%%%%%%%%%%%%%%%%%%%%%%%%%%%%%%%%%%%%%%%%%%%%%%%%%%%%%%%%%%%%%%%%%%%%%%%%%%%%%%%%%%%%%%%%
      \begin{block}{Available Compilers}
        \vspace*{-2ex}
      \begin{center}
      \begin{tabular}{|c|c|c|c|}
      \hline 
      \textbf{Language} & \textbf{Intel} & \textbf{GNU} & \textbf{PGI} \\ 
      \hline 
      Fortran77 & ifort & g77 & pgf77 \\ 
      \hline 
      Fortran90 & ifort & gfortran & pgf90 \\ 
      \hline 
      Fortran95 & ifort & gfortran & pgf95 \\ 
      \hline 
      C & icc & gcc & pgcc \\ 
      \hline 
      C++ & icpc & g++ & pgCC \\ 
      \hline 
      Debug & idb & gdb & pgdbg \\ 
      \hline
      Profile & vtune & gprof & gpprof \\ 
      \hline 
      \end{tabular} 
      \end{center}
        \vspace*{-2ex}
      \end{block}
      %%%%%%%%%%%%%%%%%%%%%%%%%%%%%%%%%%%%%%%%%%%%%%%%%%%%%%%%%%%%%%%%%%%%%%%%%%%%%%%%%%%%%%%%%%%%%%%%%%%%%%%%%%%%
      \end{column}
     \hspace*{-0.3cm}
     \begin{column}{.35\linewidth}
      %%%%%%%%%%%%%%%%%%%%%%%%%%%%%%%%%%%%%%%%%%%%%%%%%%%%%%%%%%%%%%%%%%%%%%%%%%%%%%%%%%%%%%%%%%%%%%%%%%%%%%%%%%%%
      \begin{block}{Available MPIs}
        \vspace*{-2ex}
      \begin{center}
      \begin{tabular}{|c|c|}
      \hline 
      \textbf{MPI} & \textbf{version} \\ 
      \hline 
      Intel MPI & 4.1.0,4.1.1  \\ 
      \hline 
      OpenMPI & 1.4,1.6,1.8  \\ 
      \hline 
      MPICH2 & 1.5,3.0.4  \\ 
      \hline 
      PlatformMPI & 08.02  \\ 
      \hline 
      MVAPICH2 & 1.4.1p1  \\ 
      \hline 
      \end{tabular} 
      \end{center}
        \vspace*{-2ex}
      \end{block}
      %%%%%%%%%%%%%%%%%%%%%%%%%%%%%%%%%%%%%%%%%%%%%%%%%%%%%%%%%%%%%%%%%%%%%%%%%%%%%%%%%%%%%%%%%%%%%%%%%%%%%%%%%%%%
      \end{column}
  \end{columns}

 
       %%%%%%%%%%%%%%%%%%%%%%%%%%%%%%%%%%%%%%%%%%%%%%%%%%%%%%%%%%%%%%%%%%%%%%%%%%%%%%%%%%%%%%%%%%%%%%%%%%%%%%%%%%%%
      \begin{block}{Optimization flags}
        \vspace*{-2ex}
           \begin{center}
           \begin{tabular}{|c|c|c|c|}
           \hline 
           \textbf{Compiler} & \textbf{Intel} & \textbf{GNU} & \textbf{PGI} \\ 
           \hline 
           High Opt. & -fast & -O3 -ffast-math & -fast -Mipa=fast,inline \\ 
           \hline 
           OpenMP & -openmp & -fopenmp & -mp=nonuma \\ 
           \hline 
           Debug & -g & -g & -g \\ 
           \hline 
           Profile & -p & -pg & -p \\ 
           \hline 
           Westmere & -mtarget & -march=corei7  & -tp=nehalem-64 \\ 
           \hline 
           SandyBridge & -mtarget & -march=corei7-avx  & -tp=sandybridge-64 \\ 
           \hline 
           SSE & -xsse4.2 & -msse4.2 & -Mvect=[prefetch,sse] \\ 
           \hline 
           AVX$^1$ & -xavx & -mavx & -fast \\ 
           \hline 
           \end{tabular} 
           \begin{footnotesize}
            \\1. Advanced Vector Extension (AVX) streaming SIMD instructions. Sandy Bridge processor only.
            %2. for SandyBridge use : -march=corei7-avx\\
           \end{footnotesize}
           \end{center}
        \vspace*{-2ex}
      \end{block}
      %%%%%%%%%%%%%%%%%%%%%%%%%%%%%%%%%%%%%%%%%%%%%%%%%%%%%%%%%%%%%%%%%%%%%%%%%%%%%%%%%%%%%%%%%%%%%%%%%%%%%%%%%%%%
     
        %%%%%%%%%%%%%%%%%%%%%%%%%%%%%%%%%%%%%%%%%%%%%%%%%%%%%%%%%%%%%%%%%%%%%%%%%%%%%%%%%%%%%%%%%%%%%%%%%%%%%%%%%%%%
      \begin{block}{Link MKL with OpenMPI, CDFT, ScaLAPACK, BLACS and Intel Compilers}
      \textbf{Link line:} \verb|-L${MKLROOT}/lib/intel64 -lmkl_scalapack_ilp64 \|
      \verb|        -lmkl_cdft_core -lmkl_intel_ilp64 -lmkl_sequential \|
      \verb|        -lmkl_core -lmkl_blacs_intelmpi_ilp64 -lpthread -lm|\\
      \textbf{Compiler options: }      \verb|-DMKL_ILP64 -I${MKLROOT}/include|\\
      More information at \url{http://software.intel.com/sites/products/mkl/}
      \end{block}
      %%%%%%%%%%%%%%%%%%%%%%%%%%%%%%%%%%%%%%%%%%%%%%%%%%%%%%%%%%%%%%%%%%%%%%%%%%%%%%%%%%%%%%%%%%%%%%%%%%%%%%%%%%%%
 
 
         %%%%%%%%%%%%%%%%%%%%%%%%%%%%%%%%%%%%%%%%%%%%%%%%%%%%%%%%%%%%%%%%%%%%%%%%%%%%%%%%%%%%%%%%%%%%%%%%%%%%%%%%%%%%
      \begin{block}{Link MKL with OpenMP and Intel compilers}
      \textbf{Link line:} \verb|-L${MKLROOT}/lib/intel64 -lmkl_intel_ilp64 \|\\
      \verb|       -lmkl_intel_thread -lmkl_core -lpthread -lm|\\
      \textbf{Compiler options: } \verb|-openmp -DMKL_ILP64 -I${MKLROOT}/include|
      \end{block}
      %%%%%%%%%%%%%%%%%%%%%%%%%%%%%%%%%%%%%%%%%%%%%%%%%%%%%%%%%%%%%%%%%%%%%%%%%%%%%%%%%%%%%%%%%%%%%%%%%%%%%%%%%%%%
 
 
       %%%%%%%%%%%%%%%%%%%%%%%%%%%%%%%%%%%%%%%%%%%%%%%%%%%%%%%%%%%%%%%%%%%%%%%%%%%%%%%%%%%%%%%%%%%%%%%%%%%%%%%%%%%%
      \begin{block}{Get Involved!}
       \centering{\includegraphics[scale=0.6]{img/avatar_19476.png}\\
       {\Large We would like to encourage you to send suggestions and feedback to the NeSI Team (\url{support@nesi.org.nz}).}}
      \end{block}
      %%%%%%%%%%%%%%%%%%%%%%%%%%%%%%%%%%%%%%%%%%%%%%%%%%%%%%%%%%%%%%%%%%%%%%%%%%%%%%%%%%%%%%%%%%%%%%%%%%%%%%%%%%%%
 
     \end{column}

%%%%%%%%%%%%%%%%%%%%%%%%%%%%%%%%%%%%%%%%%%%%%%%%%%%%%%%%%%%%%%
%%%%%%%%%%%%%%%%%%%%%%%%%%%%%%%%%%%%%%%%%%%%%%%%%%%%%%%%%%%%%%
%%%%%%%%%%%%%%%%%%%%%%%%%%%%%%%%%%%%%%%%%%%%%%%%%%%%%%%%%%%%%%
    
    \begin{column}{.32\linewidth}
    \vskip-2ex
%    \begin{columns}[t]
%    \begin{column}{.38\linewidth}
      %%%%%%%%%%%%%%%%%%%%%%%%%%%%%%%%%%%%%%%%%%%%%%%%%%%%%%%%%%%%%%%%%%%%%%%%%%%%%%%%%%%%%%%%%%%%%%%%%%%%%%%%%%%%
      \begin{block}{SSH Access : Login Node}
      The login node is not for running jobs, it is only for file management and job submission.\\
      \textbf{Parameters}
        \begin{itemize}
        \item host: login-01.uoa.nesi.org.nz
        \item port: 22
        \end{itemize}
      \end{block}
      %%%%%%%%%%%%%%%%%%%%%%%%%%%%%%%%%%%%%%%%%%%%%%%%%%%%%%%%%%%%%%%%%%%%%%%%%%%%%%%%%%%%%%%%%%%%%%%%%%%%%%%%%%%%
%      \end{column}
%     \hspace*{-0.6cm}
%     \begin{column}{.56\linewidth}
     
%      \end{column}
%  \end{columns}

      %%%%%%%%%%%%%%%%%%%%%%%%%%%%%%%%%%%%%%%%%%%%%%%%%%%%%%%%%%%%%%%%%%%%%%%%%%%%%%%%%%%%%%%%%%%%%%%%%%%%%%%%%%%%
      \begin{block}{Suggested Software}
        \begin{itemize}
        \item \textbf{mobaxterm} (Windows) : \url{http://mobaxterm.mobatek.net} 
        \item \textbf{Putty} (Windows) : \url{http://www.chiark.greenend.org.uk/~sgtatham/putty/}
        \item \textbf{Terminal} (MacOSX) :  (Included in the OS)
        \item \textbf{iTerm2} (MacOSX) : \url{http://www.iterm2.com}
        \item \textbf{Konsole} (Linux) : \url{http://konsole.kde.org}
        \item \textbf{GnomeTerminal} (Linux) : \url{https://wiki.gnome.org/Apps/Terminal}
        \item \textbf{yakuake} (Linux) : \url{http://yakuake.kde.org}
        \end{itemize}
      \end{block}
      %%%%%%%%%%%%%%%%%%%%%%%%%%%%%%%%%%%%%%%%%%%%%%%%%%%%%%%%%%%%%%%%%%%%%%%%%%%%%%%%%%%%%%%%%%%%%%%%%%%%%%%%%%%%
      

      %%%%%%%%%%%%%%%%%%%%%%%%%%%%%%%%%%%%%%%%%%%%%%%%%%%%%%%%%%%%%%%%%%%%%%%%%%%%%%%%%%%%%%%%%%%%%%%%%%%%%%%%%%%%
      \begin{block}{Remote File System Access}
       In order to access the file system (/home) remotely from your machine, we recommend:
        \begin{itemize}
        \item \textbf{SSHFS} (MacOSX) : \url{http://code.google.com/p/macfuse/}
        \item \textbf{SSHFS} (Linux) : \url{http://fuse.sourceforge.net/sshfs.html}
        %\item \textbf{SSHFS} (Windows) : \url{http://code.google.com/p/win-sshfs/}
        \item \textbf{Konqueror} (KDE) : type fish://user@host:port
        \item \textbf{Nautilus} (Gnome) : type sftp://user@host:port
        \item \textbf{WinSCP} (Windows) : \url{http://winscp.net}
        \end{itemize}
      \end{block}
      %%%%%%%%%%%%%%%%%%%%%%%%%%%%%%%%%%%%%%%%%%%%%%%%%%%%%%%%%%%%%%%%%%%%%%%%%%%%%%%%%%%%%%%%%%%%%%%%%%%%%%%%%%%%
 
       %%%%%%%%%%%%%%%%%%%%%%%%%%%%%%%%%%%%%%%%%%%%%%%%%%%%%%%%%%%%%%%%%%%%%%%%%%%%%%%%%%%%%%%%%%%%%%%%%%%%%%%%%%%%
      \begin{block}{Remote File System Transfer with RSYNC (Unix Only)}
       \textbf{RSYNC} over SSH protocol is the best choice to transfer big data volumes.
        \begin{itemize}
        \item Transfer data from your machine to the server:\\ \verb|rsync -avHl  /path/origin/* sshserver:/path/destination/|
        \item Transfer data from the server to your machine:\\ \verb|rsync -avHl  sshserver:/path/destination/* /path/origin/|
        \end{itemize}
      \end{block}
      %%%%%%%%%%%%%%%%%%%%%%%%%%%%%%%%%%%%%%%%%%%%%%%%%%%%%%%%%%%%%%%%%%%%%%%%%%%%%%%%%%%%%%%%%%%%%%%%%%%%%%%%%%%%

       %%%%%%%%%%%%%%%%%%%%%%%%%%%%%%%%%%%%%%%%%%%%%%%%%%%%%%%%%%%%%%%%%%%%%%%%%%%%%%%%%%%%%%%%%%%%%%%%%%%%%%%%%%%%
      \begin{block}{Remote File System Transfer with scp/sftp (Unix Only)}
       \textbf{SCP} and \textbf{SFTP} are the most popular software to transfer data across the SSH protocol.
        \begin{itemize}
        \item \textbf{SCP}: \verb|scp -pr sshserver:/path/destination/* path/destination/|
        \item \textbf{SFTP}: \verb|sftp sshserver:/path/destination/* path/destination/|
        \end{itemize}
      \end{block}
      %%%%%%%%%%%%%%%%%%%%%%%%%%%%%%%%%%%%%%%%%%%%%%%%%%%%%%%%%%%%%%%%%%%%%%%%%%%%%%%%%%%%%%%%%%%%%%%%%%%%%%%%%%%%

    \vskip-2ex
    \begin{columns}[t]
     \hspace*{+0.01cm}
    \begin{column}{.48\linewidth}
      %%%%%%%%%%%%%%%%%%%%%%%%%%%%%%%%%%%%%%%%%%%%%%%%%%%%%%%%%%%%%%%%%%%%%%%%%%%%%%%%%%%%%%%%%%%%%%%%%%%%%%%%%%%%
   %%%%%%%%%%%%%%%%%%%%%%%%%%%%%%%%%%%%%%%%%%%%%%%%%%%%%%%%%%%%%%%%%%%%%%%%%%%%%%%%%%%%%%%%%%%%%%%%%%%%%%%%%%%%
      \begin{block}{Bash Shortcuts}
      \begin{itemize}
         \item Ctrl+c - halts the current command.
         \item Ctrl+z - stops the current command.
         \item Ctrl+d - log out of current session.
         \item Ctrl+k - Delete from the cursor to the end of the line.
         \item Ctrl+w - Delete from the cursor to the start of the word.
         \item Ctrl+u - Delete from the cursor to the beginning of the line.
         \item Ctrl+r - reverse search to bash history.
         \item !! - repeats the last command.
        \end{itemize}
      \end{block}
      %%%%%%%%%%%%%%%%%%%%%%%%%%%%%%%%%%%%%%%%%%%%%%%%%%%%%%%%%%%%%%%%%%%%%%%%%%%%%%%%%%%%%%%%%%%%%%%%%%%%%%%%%%%%

      %%%%%%%%%%%%%%%%%%%%%%%%%%%%%%%%%%%%%%%%%%%%%%%%%%%%%%%%%%%%%%%%%%%%%%%%%%%%%%%%%%%%%%%%%%%%%%%%%%%%%%%%%%%%
      \end{column}
     \hspace*{-0.8cm}
     \begin{column}{.49\linewidth}
     %%%%%%%%%%%%%%%%%%%%%%%%%%%%%%%%%%%%%%%%%%%%%%%%%%%%%%%%%%%%%%%%%%%%%%%%%%%%%%%%%%%%%%%%%%%%%%%%%%%%%%%%%%%%
      \begin{block}{Process management}
      \begin{itemize}
         \item ps - display all active processes.
         \item top - display all running processes.
         \item kill pid - kill process id pid.
         \item killall proc - kill all processes named proc.
         \item bg - lists stopped or background jobs.
         \item bg n - resume job n in the background.
         \item fg - brings the most recent job to foreground.
         \item fg n - brings job n to the foreground.
        \end{itemize}
      \end{block}
      %%%%%%%%%%%%%%%%%%%%%%%%%%%%%%%%%%%%%%%%%%%%%%%%%%%%%%%%%%%%%%%%%%%%%%%%%%%%%%%%%%%%%%%%%%%%%%%%%%%%%%%%%%%%
      \end{column}
  \end{columns}
   

    \end{column}
  \end{columns}
\end{frame}
%%%%%%%%%%%%%%%%%%%%%%%%%%%%%%%%%%%%%%%%%%%%%%%%%%%%%%%%%%%%%%
%                                                        NEW FRAME                                                     %
%%%%%%%%%%%%%%%%%%%%%%%%%%%%%%%%%%%%%%%%%%%%%%%%%%%%%%%%%%%%%%
 \begin{frame}[fragile]{} 
  \begin{columns}[t]
    \begin{column}{.35\linewidth}


  %%%%%%%%%%%%%%%%%%%%%%%%%%%%%%%%%%%%%%%%%%%%%%%%%%%%%%%%%%%%%%%%%%%%%%%%%%%%%%%%%%%%%%%%%%%%%%%%%%%%%%%%%%%%
      \begin{block}{Slurm Commands}
        \begin{itemize}
        \item  \textbf{sbatch} - submits a script job.
        \item \textbf{scancel} - cancels a running or pending job.
	    \item \textbf{sinfo} - provides information on partitions and nodes.
	    \item \textbf{sview} - GUI to view job, node and partition information (ssh -X option required).
	    \item \textbf{smap} - CLI to view job, node and partition information.
        \item  \textbf{squeue} - shows the status of jobs. 
         \begin{itemize}
              \item \verb|squeue -l -j <job id>| gives extended job record information.
              \item \verb|squeue -u <user name>| lists only jobs initiated by the user.
        \end{itemize}
        \item \textbf{sbcast} - transfer file to a compute nodes allocated to a job.
        \item \textbf{interactive} - opens an interactive job session.
	    \item \textbf{sattach} - connects stdin/out/err for an existing job or job step.
        \end{itemize}
       More information about Slurm in the NeSI Slurm User Guide : \url{https://wiki.auckland.ac.nz/display/CER/Slurm+User+Guide}
      \end{block}
      %%%%%%%%%%%%%%%%%%%%%%%%%%%%%%%%%%%%%%%%%%%%%%%%%%%%%%%%%%%%%%%%%%%%%%%%%%%%%%%%%%%%%%%%%%%%%%%%%%%%%%%%%%%%

     %%%%%%%%%%%%%%%%%%%%%%%%%%%%%%%%%%%%%%%%%%%%%%%%%%%%%%%%%%%%%%%%%%%%%%%%%%%%%%%%%%%%%%%%%%%%%%%%%%%%%%%%%%%%
      \begin{block}{Slurm Submit Script Syntax}
        \begin{itemize}
        \item \verb|-A uoa99999| User account (i.e. nesiXXXXX, landcareXXXXX, uooXXXXX,uocXXXXX)
        \item \verb|--mem-per-cpu=8132| = memory/cpu (in MB)
        \item \verb|-J My_Job_Name| Sets the job name
        \item \verb|-o test-%j.%N.out| name of output file (default is slurm-jobid.out)
        \item \verb|-e test-%j.%N.err| name of error file (default is slurm-jobid.err)
        \item \verb|--mail-type=ALL| Specifies under which conditions Slurm will send out email notifications to the specified address about the state of the job (ALL, BEGIN, END, FAIL, REQUEUE).
%        \begin{itemize}
%            \item ALL - notifies about every change in the job status.
%            \item BEGIN - notifies when the job is activated.
%            \item END - notifies when the job ends.
%            \item FAIL - notifies only if the job fails.
%            \item REQUEUE - notifies only if the job is requeued.
%        \end{itemize}
        \item \verb|--mail-user=username@nesi.org.nz| user email
        \item \verb|--cpus-per-task=8| Sets the number of cores to be allocated for each task.
        \item \verb|--gres=gpu:1| - Specifies the number of GPU devices requested.
        \item \verb|--gres=mic:1| - Specifies the number of Intel Xeon Phi Co-Processor devices requested.
        \item \verb|--ntasks-per-node=16| specifies the total number of tasks to be run on each available node.
        \item \verb|--ntasks=96| specifies how many tasks (cores) are to be run in total
        \item \verb|--time=01:00:00| Sets the limit for the elapsed time for which a job can run
        \item \verb|--nodes=N-M| Defines the minimum (N) and the maximum (M) number of nodes.
        \item \verb|-C sb| Requires an specific hardware architecture (sb=Sandybridge,wm=Westmere)
        \item \verb|--exclusive| Requires exclusive execution
        \end{itemize}
      \end{block}
      %%%%%%%%%%%%%%%%%%%%%%%%%%%%%%%%%%%%%%%%%%%%%%%%%%%%%%%%%%%%%%%%%%%%%%%%%%%%%%%%%%%%%%%%%%%%%%%%%%%%%%%%%%%%

     %%%%%%%%%%%%%%%%%%%%%%%%%%%%%%%%%%%%%%%%%%%%%%%%%%%%%%%%%%%%%%%%%%%%%%%%%%%%%%%%%%%%%%%%%%%%%%%%%%%%%%%%%%%%
      \begin{block}{Interactive sessions}
      The interactive sessions will allow you to build the binaries for specific architecture. The binaries compiled in Westmere can run in Sandy Bridge, but it can \textbf{NOT} exploit all the Sandy bridge features. The binaries compiled in Sandy Bridge can \textbf{NOT} run in the Westmere nodes. The interactively usage is limited up to \textbf{24h} of \textbf{walltime}.
      \vspace*{-0.2cm}
	 \begin{verbatim}
Usage: interactive [-A] [-a] [-c] [-m] [-J] [-e]
Mandatory arguments:
	 -A: account
Optional arguments:
	 -a: architecture (default: wm, values sb=SandyBridge wm=Westmere)
	 -c: number of CPU cores (default: 1)
	 -m: amount of memory (GB) per core (default: 1 [GB])
	 -J: job name
	 -e: binary that you want to run interactively
example : interactive -A nesi99999
	\end{verbatim} 
      \end{block}
      %%%%%%%%%%%%%%%%%%%%%%%%%%%%%%%%%%%%%%%%%%%%%%%%%%%%%%%%%%%%%%%%%%%%%%%%%%%%%%%%%%%%%%%%%%%%%%%%%%%%%%%%%%%%


    \end{column}

%%%%%%%%%%%%%%%%%%%%%%%%%%%%%%%%%%%%%%%%%%%%%%%%%%%%%%%%%%%%%%
%%%%%%%%%%%%%%%%%%%%%%%%%%%%%%%%%%%%%%%%%%%%%%%%%%%%%%%%%%%%%%
%%%%%%%%%%%%%%%%%%%%%%%%%%%%%%%%%%%%%%%%%%%%%%%%%%%%%%%%%%%%%%
    
    \begin{column}{.3\linewidth}
    

      %%%%%%%%%%%%%%%%%%%%%%%%%%%%%%%%%%%%%%%%%%%%%%%%%%%%%%%%%%%%%%%%%%%%%%%%%%%%%%%%%%%%%%%%%%%%%%%%%%%%%%%%%%%%
      \begin{block}{Submit Script Example : Serial}
              \vspace*{-3ex}
        \begin{verbatim}
#!/bin/bash
#SBATCH -J OpenMP_JOB
#SBATCH -A uoa99999         # Project Account
#SBATCH --time=01:00:00     # Walltime
#SBATCH --mem-per-cpu=8132  # memory/cpu (in MB)
srun serial_binary
        \end{verbatim}
                \vspace*{-4ex}
      \end{block}
      %%%%%%%%%%%%%%%%%%%%%%%%%%%%%%%%%%%%%%%%%%%%%%%%%%%%%%%%%%%%%%%%%%%%%%%%%%%%%%%%%%%%%%%%%%%%%%%%%%%%%%%%%%%%

      %%%%%%%%%%%%%%%%%%%%%%%%%%%%%%%%%%%%%%%%%%%%%%%%%%%%%%%%%%%%%%%%%%%%%%%%%%%%%%%%%%%%%%%%%%%%%%%%%%%%%%%%%%%%
      \begin{block}{Submit Script Example : OpenMP}
              \vspace*{-3ex}
        \begin{verbatim}
#!/bin/bash
#SBATCH -J OpenMP_JOB
#SBATCH -A uoa99999         # Project Account
#SBATCH --time=01:00:00     # Walltime
#SBATCH --mem-per-cpu=8132  # memory/cpu (in MB)
#SBATCH --cpus-per-task=8   # 8 OpenMP Threads
srun openmp_binary
        \end{verbatim}
                \vspace*{-4ex}
      \end{block}
      %%%%%%%%%%%%%%%%%%%%%%%%%%%%%%%%%%%%%%%%%%%%%%%%%%%%%%%%%%%%%%%%%%%%%%%%%%%%%%%%%%%%%%%%%%%%%%%%%%%%%%%%%%%%

      %%%%%%%%%%%%%%%%%%%%%%%%%%%%%%%%%%%%%%%%%%%%%%%%%%%%%%%%%%%%%%%%%%%%%%%%%%%%%%%%%%%%%%%%%%%%%%%%%%%%%%%%%%%%
      \begin{block}{Submit Script Example : MPI}
              \vspace*{-3ex}
        \begin{verbatim}
#!/bin/bash
#SBATCH -J MPI_JOB
#SBATCH -A uoa99999         # Project Account
#SBATCH --time=01:00:00     # Walltime
#SBATCH --ntasks=2          # number of tasks
#SBATCH --mem-per-cpu=8132  # memory/cpu (in MB)
srun mpi_binary
        \end{verbatim}
                \vspace*{-4ex}
      \end{block}
      %%%%%%%%%%%%%%%%%%%%%%%%%%%%%%%%%%%%%%%%%%%%%%%%%%%%%%%%%%%%%%%%%%%%%%%%%%%%%%%%%%%%%%%%%%%%%%%%%%%%%%%%%%%%

      %%%%%%%%%%%%%%%%%%%%%%%%%%%%%%%%%%%%%%%%%%%%%%%%%%%%%%%%%%%%%%%%%%%%%%%%%%%%%%%%%%%%%%%%%%%%%%%%%%%%%%%%%%%%
      \begin{block}{Submit Script Example : Hybrid (MPI+OpenMP)}
              \vspace*{-3ex}
        \begin{verbatim}
#!/bin/bash
#SBATCH -J Hybrid_JOB
#SBATCH -A uoa99999         # Project Account
#SBATCH --time=01:00:00     # Walltime
#SBATCH --ntasks=4          # number of tasks
#SBATCH --mem-per-cpu=8132  # memory/cpu (in MB)
#SBATCH --cpus-per-task=8   # 8 OpenMP Threads
#SBATCH --nodes=1           # number nodes
srun binary_hybrid
        \end{verbatim}
                \vspace*{-4ex}
      \end{block}
      %%%%%%%%%%%%%%%%%%%%%%%%%%%%%%%%%%%%%%%%%%%%%%%%%%%%%%%%%%%%%%%%%%%%%%%%%%%%%%%%%%%%%%%%%%%%%%%%%%%%%%%%%%%%

      %%%%%%%%%%%%%%%%%%%%%%%%%%%%%%%%%%%%%%%%%%%%%%%%%%%%%%%%%%%%%%%%%%%%%%%%%%%%%%%%%%%%%%%%%%%%%%%%%%%%%%%%%%%%
      \begin{block}{Submit Script Example : Hybrid (MPI+CUDA)}
              \vspace*{-3ex}
        \begin{verbatim}
#!/bin/bash
#SBATCH -J GPU_JOB
#SBATCH --time=01:00:00     # Walltime
#SBATCH -A uoa99999         # Project Account
#SBATCH --ntasks=4          # number of tasks
#SBATCH --ntasks-per-node=2 # number of tasks per node
#SBATCH --mem-per-cpu=8132  # memory/cpu (in MB)
#SBATCH --cpus-per-task=4   # 4 OpenMP Threads
# The following line will request GPUs per node. In this  
# particular example, it means 4 GPUs in total.
#SBATCH --gres=gpu:2       
#SBATCH -C kepler
srun binary_cuda_mpi
        \end{verbatim}
                \vspace*{-4ex}
      \end{block}
      %%%%%%%%%%%%%%%%%%%%%%%%%%%%%%%%%%%%%%%%%%%%%%%%%%%%%%%%%%%%%%%%%%%%%%%%%%%%%%%%%%%%%%%%%%%%%%%%%%%%%%%%%%%%


    \end{column}

%%%%%%%%%%%%%%%%%%%%%%%%%%%%%%%%%%%%%%%%%%%%%%%%%%%%%%%%%%%%%%
%%%%%%%%%%%%%%%%%%%%%%%%%%%%%%%%%%%%%%%%%%%%%%%%%%%%%%%%%%%%%%
%%%%%%%%%%%%%%%%%%%%%%%%%%%%%%%%%%%%%%%%%%%%%%%%%%%%%%%%%%%%%%
    
    \begin{column}{.35\linewidth}
    
      %%%%%%%%%%%%%%%%%%%%%%%%%%%%%%%%%%%%%%%%%%%%%%%%%%%%%%%%%%%%%%%%%%%%%%%%%%%%%%%%%%%%%%%%%%%%%%%%%%%%%%%%%%%%
      \begin{block}{Submit Script Example : Job Array}
              \vspace*{-3ex}
        \begin{verbatim}
#!/bin/bash
#SBATCH -J JobArray
#SBATCH --time=01:00:00     # Walltime
#SBATCH -A uoa99999         # Project Account
#SBATCH --ntasks=1          # number of tasks
#SBATCH --mem-per-cpu=8132  # memory/cpu (in MB)
#SBATCH --cpus-per-task=4   # 4 OpenMP Threads
#SBATCH --array=1-1000      # Array definition
srun binary_array $SLURM_ARRAY_TASK_ID
        \end{verbatim}
                \vspace*{-4ex}
      \end{block}
      %%%%%%%%%%%%%%%%%%%%%%%%%%%%%%%%%%%%%%%%%%%%%%%%%%%%%%%%%%%%%%%%%%%%%%%%%%%%%%%%%%%%%%%%%%%%%%%%%%%%%%%%%%%%


        %%%%%%%%%%%%%%%%%%%%%%%%%%%%%%%%%%%%%%%%%%%%%%%%%%%%%%%%%%%%%%%%%%%%%%%%%%%%%%%%%%%%%%%%%%%%%%%%%%%%%%%%%%%%
      \begin{block}{User Environment}
      \textbf{LMOD} is very useful to manage environment variables for each application and it is very easy to use. It loads the needed environment by a certain application and its dependencies automatically. The command line is fully compatible with the previous \textbf{Environment Modules}, and it provides simple short-cuts and advanced features.\\
      Syntax : \verb|module [options] sub-command [args ...]|\\
      \textbf{Loading/Unloading sub-commands}
        \begin{itemize}
        \item  \textbf{load $|$ add} load module(s)
        \item  \textbf{del $|$ unload} Remove module(s), do not complain if not found
        \item  \textbf{purge} unload all modules
        \item  \textbf{update} reload all currently loaded modules.
        \end{itemize}
       \textbf{Listing / Searching sub-commands}
        \begin{itemize}
        \item  \textbf{list} List loaded modules
        \item  \textbf{avail $|$ av} List available modules
        \item  \textbf{avail $|$ av string} List available modules that contain "string".
        \item  \textbf{spider} List all possible modules
        \item  \textbf{spider module} List all possible version of that module file
        \item  \textbf{spiderstring} List all module that contain the "string".
        \end{itemize}
       \textbf{Short-cuts}
       \begin{itemize}
        \item \textbf{ml} - means: module list
        \item \textbf{ml foo bar} - means: module load foo bar
        \item \textbf{ml -foo -bar baz goo}- means: module unload foo bar; module load baz goo;
        \end{itemize}
        More information at \url{http://www.tacc.utexas.edu/tacc-projects/lmod}
      \end{block}
      %%%%%%%%%%%%%%%%%%%%%%%%%%%%%%%%%%%%%%%%%%%%%%%%%%%%%%%%%%%%%%%%%%%%%%%%%%%%%%%%%%%%%%%%%%%%%%%%%%%%%%%%%%%%
      
    \end{column}
  \end{columns}
\end{frame}

\end{document}


%%%%%%%%%%%%%%%%%%%%%%%%%%%%%%%%%%%%%%%%%%%%%%%%%%%%%%%%%%%%%%%%%%%%%%%%%%%%%%%%%%%%%%%%%%%%%%%%%%%%
%%% Local Variables: 
%%% mode: latex
%%% TeX-PDF-mode: t
%%% End: 
